\documentclass[11pt]{article}
\usepackage[margin=1in]{geometry}
\usepackage{amsmath,amssymb}

\title{MATH 375: Probability Theory}
\author{Mariana Restrepo}
\date{Fall 2026}

\begin{document}
\maketitle

\begin{center}
\large Lecture 3
\end{center}

\tableofcontents

\section{Lecture 3: Events and Counting Probabilities}

\subsection{Reading event notation}

Let $E$ and $F$ be events. Then:
\[
\begin{array}{ll}
E\cup F & \text{$E$ occurs, $F$ occurs, or both occur},\\[2pt]
E\cap F & \text{both $E$ and $F$ occur},\\[2pt]
E^c & \text{$E$ does not occur}.
\end{array}
\]

Two useful probability rules are
\[
P(E^c)=1-P(E)
\]
and
\[
P(E\cup F)=P(E)+P(F)-P(E\cap F).
\]
The intersection is subtracted because it was counted once in
$P(E)$ and again in $P(F)$.

We can also split $F$ into two nonoverlapping pieces:
\[
F=(F\cap E)\cup(F\cap E^c).
\]
Consequently,
\[
P(F)=P(F\cap E)+P(F\cap E^c).
\]

\subsection{When one event is contained in another}

We write $E\subseteq F$ when every outcome in $E$ is also in $F$.
For example, when rolling a die, the event ``roll a 6'' is contained
in the event ``roll an even number.''

If $E\subseteq F$, then $E\cap F=E$. Using the previous rule,
\[
\begin{aligned}
P(F)
&=P(F\cap E)+P(F\cap E^c)\\
&=P(E)+P(F\cap E^c)\\
&\geq P(E).
\end{aligned}
\]
Therefore,
\[
\boxed{E\subseteq F\quad\Longrightarrow\quad P(E)\leq P(F).}
\]

\subsection{Example: liking neither book}

Jane likes the first book with probability $0.5$, the second with
probability $0.4$, and both with probability $0.3$.

Let $A$ mean ``likes the first book'' and $B$ mean ``likes the
second book.'' The event ``likes neither'' is $(A\cup B)^c$.
First find
\[
P(A\cup B)=P(A)+P(B)-P(A\cap B)
=0.5+0.4-0.3=0.6.
\]
Then use the complement:
\[
\boxed{P\bigl((A\cup B)^c\bigr)=1-0.6=0.4.}
\]

\subsection{Equally likely outcomes}

If every outcome in a finite sample space $S$ is equally likely,
then for an event $E$,
\[
P(E)=\frac{|E|}{|S|}.
\]

\paragraph{Example: two dice.}
Two distinguishable fair dice have $6\cdot6=36$ ordered outcomes.
The outcomes whose sum is 6 are
\[
(1,5),\ (2,4),\ (3,3),\ (4,2),\ (5,1).
\]
There are five favorable outcomes, so
\[
\boxed{P(\text{sum is 6})=\frac{5}{36}.}
\]

\subsection{Example: selecting a committee}

Select five people uniformly from a group of six men and nine women.
What is the probability of selecting exactly three men and two women?

There are $\binom{15}{5}$ possible committees. For a favorable
committee, choose three of the six men and two of the nine women:
\[
|E|=\binom{6}{3}\binom{9}{2}.
\]
Thus,
\[
\boxed{
P(E)=\frac{\binom{6}{3}\binom{9}{2}}{\binom{15}{5}}
=\frac{240}{1001}.
}
\]

\subsection{Example: a five-card straight}

A five-card hand is a \textbf{straight} if its ranks are consecutive
and its cards are not all of the same suit. Aces may be low
($A,2,3,4,5$) or high ($10,J,Q,K,A$), giving ten possible rank
sequences.

For each sequence, there are $4^5$ choices of suits. Exclude the
four choices where all five cards have the same suit. Therefore,
\[
|E|=10(4^5-4).
\]
Since there are $\binom{52}{5}$ five-card hands,
\[
\boxed{
P(\text{straight})
=\frac{10(4^5-4)}{\binom{52}{5}}.
}
\]

\subsection{Example: bridge hands}

In bridge, four labeled players each receive 13 cards.

\paragraph{One specified player receives all 13 spades.}
There are $\binom{52}{13}$ possible hands for that player, and only
one contains all the spades. Hence
\[
P(\text{player 1 gets all spades})
=\frac{1}{\binom{52}{13}}.
\]

\paragraph{Any one player receives all 13 spades.}
The four possible recipients give mutually exclusive events, so
\[
\boxed{
P(\text{some player gets all spades})
=\frac{4}{\binom{52}{13}}.
}
\]

\paragraph{Each player receives exactly one ace.}
The number of ways to deal 52 cards to four labeled players is
\[
|S|=\frac{52!}{(13!)^4}.
\]

First assign the four distinct aces, one to each player, in $4!$
ways. Then divide the 48 non-ace cards into four groups of 12:
\[
|E|=4!\frac{48!}{(12!)^4}.
\]
Therefore,
\[
\boxed{
P(\text{one ace per player})
=\frac{4!\,48!/(12!)^4}{52!/(13!)^4}.
}
\]

\end{document}